\chapter{Automatic checks: continuous integration}
\label{ch:ci}

\chaptergoals{5 of 6 in Part~I.}
{\cref{ch:remotes} (pushing to GitHub) and \cref{ch:pixi} (tasks and
\cmd{-{}-locked}); your \file{pixi-playground}; a GitHub account.}
{read a GitHub Actions workflow, set one up, find out why a check failed and
reproduce it locally, and recognise the checks every ShipSoft pull request
must pass.}

\section{What CI does}

\emph{Continuous integration} (CI) means: every time somebody pushes, a fresh
machine checks out the commit, builds it, and runs the tests and the
linters. The result appears on the pull request as a \emph{check}, green or
red. Nobody has to remember to run the tests, the reviewer does not have to
run them either, and they run in a clean environment rather than on somebody's
carefully configured laptop. A repository can make checks \emph{required}: a
pull request with a red required check cannot be merged.

CI only works if it runs exactly what you run. That is what the pixi tasks and
the lock file of \cref{ch:pixi} are for: CI runs
\cmd{pixi run -{}-locked <task>}, the same command, in the same environment,
that you run at home.

\section{GitHub Actions in five words}

GitHub's CI is called GitHub Actions. Its vocabulary:

\begin{description}[style=nextline]
  \item[workflow] a YAML file in \file{.github/workflows/} in the repository.
  \item[events] what starts it, under \texttt{on:}: a push, a pull request,
    the merge queue, a schedule, a button.
  \item[jobs] what it does. Each job runs on a fresh virtual machine
    (\texttt{runs-on: ubuntu-latest}); jobs run in parallel unless one
    \texttt{needs} another.
  \item[steps] the commands of a job, in order. A step either \texttt{uses}
    a ready-made \emph{action} (check out the code, install pixi) or
    \texttt{run}s shell commands.
  \item[logs] everything a step prints, kept on GitHub, one click away from
    the pull request.
\end{description}

\section{Your first workflow}

\codefile[basicstyle=\ttfamily\footnotesize]{yaml}{code/ci/playground-ci.yml}

Three steps: \texttt{actions/checkout} fetches the commit being tested;
\texttt{prefix-dev/setup-pixi} installs pixi and the environment, from
\file{pixi.lock} (\texttt{locked: true}); and the last step runs your task.

\begin{exercise}{CI for your playground}
\label{ex:firstci}
On GitHub, create a new, empty repository \texttt{pixi-playground} (no
README, no licence). Then, in your \file{pixi-playground}:
\begin{shell}
mkdir -p .github/workflows
cp $CODE/ci/playground-ci.yml .github/workflows/ci.yml
git add .github && git commit -m "ci: run hello on every push"
git remote add origin git@github.com:<you>/pixi-playground.git
git push -u origin main
\end{shell}
Open the repository's \emph{Actions} tab and follow the run. Click on the job
and open each step: where does it print the ROOT version?
\end{exercise}

\begin{pitfall}
CI runs on Linux. \cmd{pixi init} on a Mac writes only your Mac's platform
into \file{pixi.toml}, and \cmd{pixi run -{}-locked} then fails in CI because
the lock file has nothing for \texttt{linux-64}. Add it with
\cmd{pixi workspace platform add linux-64} and commit the new lock file.
\end{pitfall}

\begin{exercise}{Red, then green}
\label{ex:redgreen}
Break the task on a branch and propose the breakage to yourself:
\begin{shell}
git switch -c broken
echo 'raise SystemExit("deliberately broken")' >> hello.py
git commit -am "test: break hello on purpose"
git push -u origin broken
\end{shell}
Open a pull request from \texttt{broken} into \texttt{main} of your own
repository. Watch the check turn red, open its log, and find the error.
Reproduce it locally with the very same command, \cmd{pixi run -{}-locked
hello}. Then repair \file{hello.py}, commit, \cmd{git push}, and watch the
same pull request turn green. Close it without merging.
\end{exercise}

\section{The workflow of the team repository}
\label{sec:teamci}

The team repository of \cref{ch:team} has a workflow laid out like the
ShipSoft ones. Read it once now; you will see it run on your pull request.

\codefile[basicstyle=\ttfamily\footnotesize]{yaml}{code/teamwork/template/.github/workflows/ci.yml}

\begin{description}[style=nextline]
  \item[Events.] Pull requests and pushes to \texttt{main}, as before, plus
    \texttt{merge\_group}: the merge queue runs the checks once more on the
    rebased commit before it lands (\cref{sec:pullrequests}).
  \item[Concurrency.] A new push to a pull request cancels the run for the
    previous push, which nobody needs any more.
  \item[Two jobs] run in parallel: \texttt{tasks} runs
    \cmd{pixi run -{}-locked check} (all ten task checks of \cref{ch:team}), and
    \texttt{lint} runs \texttt{ruff} in its own pixi environment.
  \item[One summary job,] ``All checks passed'', needs both and fails if
    either failed or was cancelled. Its name is the one required check of the
    repository. Jobs can be added or renamed later without touching the
    repository settings; \texttt{if: always()} makes it report a failure
    instead of being silently skipped.
\end{description}

\begin{exercise}{Predict the checks}
Without running anything, answer from the workflow above:
\begin{enumerate}
  \item A pull request adds a correct task~4 but also an unused
        \texttt{import}. Which jobs are green, which red, and can the pull
        request be merged?
  \item You push twice within a minute. How many runs complete?
  \item Which single command reproduces the \texttt{tasks} job on your
        machine, and which the \texttt{lint} job?
\end{enumerate}
\end{exercise}

\begin{solution}
(1) \texttt{tasks} green, \texttt{lint} red (ruff rule F401), so ``All checks
passed'' is red and the pull request cannot be merged until the import is
removed. (2) One: the second push cancels the first run
(\texttt{concurrency}). (3) \cmd{pixi run -{}-locked check} and
\cmd{pixi run -{}-locked lint}.
\end{solution}

\section{CI in the ShipSoft repositories}

The ShipSoft repositories share their CI through \emph{reusable workflows} in
\repo{.github}: each repository's workflow is a few lines that call a common
one. From \repo{Shannon}'s \file{.github/workflows/pixi-build.yml}:
\begin{lstlisting}[language=yaml]
jobs:
  build:
    uses: ShipSoft/.github/.github/workflows/pixi-cmake-build.yml@main
    with:
      tasks: '["configure", "build", "install", "test", "smoke"]'
\end{lstlisting}
followed by the same ``All checks passed'' summary job as in the team
repository. What runs, and what it means for you:

\begin{center}\small
\begin{tabularx}{\linewidth}{l L}
\toprule
Workflow & What it checks \\
\midrule
\texttt{pixi-cmake-build} & runs the listed pixi tasks with \cmd{-{}-locked}, compiling through
                            a cache (``All checks passed'') \\
\texttt{prek}             & the lint hooks of the \texttt{lint} environment on all files, and posts
                            their output on the pull request (``Lint passed'') \\
\texttt{commit-check}     & every commit message of the pull request against Conventional Commits
                            (part of ``Lint passed'') \\
\texttt{clang-tidy}       & C++ guidelines, reported only on the lines your pull request changes \\
\texttt{pixi-lock-update} & once a month, opens a pull request with an updated \file{pixi.lock} \\
\texttt{physics-metrics}  & physics distributions compared with a reference (\repo{ship-ci-metrics}) \\
\bottomrule
\end{tabularx}
\end{center}

\begin{note}
Run the lint checks before you push and the ``Lint passed'' check is a
formality: \cmd{pixi run lint} checks everything, and \cmd{pixi run
install-hooks} makes the same checks run at every \cmd{git commit}. When a
check is red anyway, the log of the failed step names the pixi task; rerun
exactly that task locally.
\end{note}

\begin{exercise}[colback=white]{Optional: read a real run}
Open \url{https://github.com/ShipSoft/aegir/actions}, choose a recent run of
``Pixi Build'', and find the log group of \cmd{pixi run -{}-locked test}. How
long did the build take, and how many tests ran?
\end{exercise}

\begin{levelcheck}
\begin{itemize}
  \can name the parts of a workflow: events, jobs, steps;
  \can set up CI for a pixi project and watch it run;
  \can go from a red check to its log, and reproduce the failure with the
       same pixi task locally;
  \can explain why a summary job, rather than every job, is the required
       check, and what \texttt{merge\_group} is for.
\end{itemize}
\end{levelcheck}

\begin{solution}
Both workflows in this chapter pass \texttt{actionlint}~1.7.12, and use the
action versions of the ShipSoft reusable workflows on 15 October 2026
(\texttt{actions/checkout@v7}, \texttt{prefix-dev/setup-pixi@v0.11.0}). They
have not been run on GitHub: do \cref{ex:firstci} once before the session.
\end{solution}
